\section{当前可达重构模型}

本章只描述 \srcpath{src/network\_reconfiguration/topology\_reconfiguration.cpp:run\_topology\_reconfiguration} 的
\srcpath{run\_topology\_reconfiguration}。历史入口
\srcpath{solve\_optimal\_reconfiguration} 在
\srcpath{src/network\_reconfiguration/topology\_analysis.cpp:solve\_optimal\_reconfiguration} 已提前包装并返回该混合求解器结果；其后保留的旧 AC MILP
在当前控制流不可达，故不作为现行模型。

\subsection{变量布局和边界传播}

统一边序为 AC、DC、VSC。拓扑变量包括有符号虚拟流 $F_e\in[-n_b,n_b]$、
源虚拟注入 $F_g\in[0,n_b]$、支路状态 $\beta_e$ 和根指示 $\gamma_g$。
启用 PF 时再创建线路/VSC 的 $P,Q$、源 $P_g,Q_g$、平方电压 $v_i$、有功切负荷
$s_i^P\ge0$、AC 无功切负荷 $s_i^Q\ge0$；仅在同时启用 loss aware 时创建
$t_e\ge0$。准确偏移见 \srcpath{src/network\_reconfiguration/topology\_reconfiguration.cpp:VarLayout}。

状态边界按以下分支执行：不参与拓扑的 overlay 边固定为原状态且虚拟流固定零；故障边
$\zeta_e<0.5$ 固定 $\beta_e=F_e=0$；合法可切换边保留 $0\le\beta_e\le1$ 并成为二元变量；
其余边固定为原状态，固定断开时同时固定 $F_e=0$。
同一当前连通分量内部的断开 tie 即便此前自由，也会固定断开。
依据为 \srcpath{src/network\_reconfiguration/topology\_reconfiguration.cpp:add\_dc\_source（lambda）}。

\subsection{实际目标函数}

对初始断开边，$\beta_e$ 的动作系数为 $+\lambda_{sw}c_e$；对初始闭合边为
$-\lambda_{sw}c_e$，被省略的常数是 $\lambda_{sw}c_e$。因此目标中的动作部分等价于
对状态变化计费，但存入 LP 的不是 $\lambda_{sw}\sum\beta_e$。
损耗分支为
\begin{equation}
 L_e=\begin{cases}
 \lambda_{loss}|r_e|t_e,&\fld{loss\_aware}\land\fld{enable\_pf},\\
 \lambda_{loss}|r_e|\beta_e,&\text{否则}.
 \end{cases}
\end{equation}
启用 PF 时再加 $\lambda_{shed}(\sum_i s_i^P+\sum_{i\in AC}s_i^Q)$；
除 primary root 外的每个 $\gamma_g$ 加 $\lambda_{island}\gamma_g$。
完整赋值位于 \srcpath{src/network\_reconfiguration/topology\_reconfiguration.cpp:log\_domain\_cardinality（lambda）}。

\subsection{拓扑等式与不等式}

统一树模式下，虚拟流平衡的矩阵符号为：from 端 $+F_e$、to 端 $-F_e$、源端
$-F_g$，每个母线右端为 $-1$；即
\begin{equation}
 \sum_{e:f(e)=i}F_e-\sum_{e:t(e)=i}F_e-\sum_{g:i(g)=i}F_g=-1.
\end{equation}
分域树模式把所有 VSC 虚拟流固定为零，并分别在每个 AC/DC 连通分量代表节点把右端
从 $-1$ 增加该分量母线数；见 \srcpath{src/network\_reconfiguration/topology\_reconfiguration.cpp:log\_domain\_cardinality（lambda）}。

统一模式的森林基数等式为
\begin{equation}
 \sum_{e\in\mathcal E_{participating}}\beta_e+\sum_g\gamma_g=n_b.
\end{equation}
分域模式不使用该等式，而对每个预计算 AC 分量固定
$\sum_{e\in component}\beta_e=|V_{component}|-1$；仅当不允许 DC mesh 时，对每个 DC 分量
也创建同式。VSC 不计入分域树边数。依据为
\srcpath{src/network\_reconfiguration/topology\_reconfiguration.cpp:log\_domain\_cardinality（lambda）}。

仅对自由 $\beta_e$ 创建
\begin{equation}
 -n_b\beta_e\le F_e\le n_b\beta_e.
\end{equation}
每个源创建 $F_g\le n_b\gamma_g$，并创建 $\sum_g\gamma_g\ge1$。
统一模式存在真实 DC 源时，另创建 DC 源集合上的 $\sum\gamma_g\ge1$。
依据为 \srcpath{src/network\_reconfiguration/topology\_reconfiguration.cpp:log\_domain\_cardinality（lambda）} 和
\srcpath{src/network\_reconfiguration/topology\_reconfiguration.cpp:bigm\_branch（lambda）}。

\subsection{可选 LinDistFlow 约束}

启用 PF 时，有功矩阵行采用 from $+P_e$、to $-P_e$、源 $-P_g$、切负荷 $-s_i^P$，
右端为 $-P_{d,i}$；AC 无功采用相同符号，右端 $-Q_{d,i}$。
这对应
\begin{align}
 \sum_{out}P_e-\sum_{in}P_e-\sum_gP_g-s_i^P&=-P_{d,i},\\
 \sum_{out}Q_e-\sum_{in}Q_e-\sum_gQ_g-s_i^Q&=-Q_{d,i}\quad(i\in AC).
\end{align}
VSC $P$ 同时进入 AC/DC 两端；VSC $Q$ 只进入属于 AC 的端点。
依据为 \srcpath{src/network\_reconfiguration/topology\_reconfiguration.cpp:log\_domain\_cardinality（lambda）}。

启用电压约束时，对 AC/DC 线路而非 VSC 创建
\begin{equation}
 |v_j-v_i+2r_eP_e+2x_eQ_e|\le M_e(1-\beta_e),
\end{equation}
DC 边没有 $Q$ 项。源码的支路大 M 为
\begin{equation}
 M_e=\min\!\left(\fld{big\_m\_v},
 \max\left(0.1,|\Delta v_{span}|+2(|r_e|+|x_e|)\bar S_e\right)\right),
\end{equation}
其中 $\Delta v_{span}$ 取两种端点上下界差的较大者；见
\srcpath{src/network\_reconfiguration/topology\_reconfiguration.cpp:log\_domain\_cardinality（lambda）}。

启用热限时，线路/VSC 有功均使用独立箱约束
$|P_e|\le\bar S_e\beta_e$；AC 线路与 VSC 无功使用
$|Q_e|\le\bar S_e\beta_e$。这不是圆形视在功率约束；DC 线路无 $Q$ 变量。
loss aware 分支创建 $P_e-t_e\le0$ 和 $-P_e-t_e\le0$，故 $t_e\ge|P_e|$；
注释称其为 $I^2R$ 代理，但实际目标是 $|r_e||P_e|$ 的线性代理。
依据为 \srcpath{src/network\_reconfiguration/topology\_reconfiguration.cpp:bigm\_branch（lambda）} 和
\srcpath{src/network\_reconfiguration/topology\_reconfiguration.cpp:bigm\_branch（lambda）}。

\subsection{开关预算与结果语义}

预算只计自由边，并乘设备动作成本 $c_e$：初始断开边贡献 $c_e\beta_e$，初始闭合边贡献
$c_e(1-\beta_e)$。矩阵形式把闭合边常数移到右端：
\begin{equation}
 \sum_{e:\alpha_e=0}c_e\beta_e-
 \sum_{e:\alpha_e=1}c_e\beta_e
 \le N_{max}-\sum_{e:\alpha_e=1}c_e.
\end{equation}
仅当 \fld{max\_switch\_ops}>0 时创建，见
\srcpath{topology\_reconfiguration.cpp:run\_topology\_reconfiguration}。

旧手册给出的 $\sum r_eS_B$ 不是当前 loss aware 目标。结果的
\fld{feasible}/\fld{optimal}/\fld{proven\_optimal} 还受后端状态字符串规范化和 primal 可用性影响；
详细提取位于 \srcpath{topology\_reconfiguration.cpp:run\_topology\_reconfiguration}。

\subsection{覆盖边界}

本章覆盖当前可达 MILP 的变量、目标和带源码标签的 T/P/G 约束。源容量聚合、设备开断能力、
overlay 去重、伪根生成、BFS 初解、后处理开关序列均是离散工程分支，源码区间分别为
\srcpath{src/network\_reconfiguration/topology\_reconfiguration.cpp:run\_topology\_reconfiguration}、
\srcpath{src/network\_reconfiguration/topology\_reconfiguration.cpp:bigm\_branch（lambda）}；本章不以通用图论公式替代这些实现。
